% chapters/ch06.tex
% chapters/ch06.tex
% 第~
\ref{chap:6}~章 \index{jinsuanzi@紧算子 (compact operator)}
紧算子. 自伴紧算子的\index{pufenjie@谱分解 (spectral decomposition)}
谱分解
% 源: 原书 PDF 第 166--188 页 (印刷页 157--179)

\chapter{紧算子. 自伴紧算子的谱分解}
\label{chap:6}

\section{定义. 初等性质. 伴随算子}
\label{sec:6.1}

本章自始至终 (除非另有说明) 用 $E$ 与 $F$ 表示两个 Banach 空间.

\begin{definition}
设 $T\in\LL(E,F)$ 为有界算子.若 $T(B_E)$ 在 $F$ 中 (在强拓扑下) 有紧的
闭包, 则称 $T$ 是\emph{紧算子}.
\end{definition}

由 $E$ 到 $F$ 的一切紧算子构成之集记作 $\KK(E,F)$.为简单起见, 记
$\KK(E)=\KK(E,E)$.

\begin{theorem}
\label{thm:6.1}
集 $\KK(E,F)$ 是 $\LL(E,F)$ 的闭线性子空间 (在由范数
$\norm{\cdot}_{\LL(E,F)}$ 所诱导的拓扑下).
\end{theorem}

\begin{proof}
显然, 两个紧算子之和是紧算子.设 $(T_n)$ 为紧算子序列, $T$ 为有界算子,
使得 $\norm{T_n-T}_{\LL(E,F)}\to0$.我们断言 $T$ 是紧算子.由于 $F$ 完备,
只须验证: 对每个 $\varepsilon>0$, 存在 $T(B_E)$ 的半径 $\varepsilon$ 的有限覆盖
(见例如 J.~R.~Munkres \bref{Munkres}{1} 第~\ref{sec:7.3}~节).取整数 $n$, 使
$\norm{T_n-T}_{\LL(E,F)}<\varepsilon/2$.由于 $T_n(B_E)$ 有紧的闭包, 存在
$T_n(B_E)$ 的半径 $\varepsilon/2$ 的有限覆盖, 即
$T_n(B_E)\subset\bigcup_{i\in I}B(f_i,\varepsilon/2)$.由此推出
\begin{equation*}
T(B_E)\subset\bigcup_{i\in I}B(f_i,\varepsilon).
\end{equation*}
\end{proof}

\begin{definition}
设 $T\in\LL(E,F)$.若 $T$ 的值域 $R(T)$ 是有限维的, 则称 $T$ 是
\emph{有限秩算子}.
\end{definition}

显然, 任何有限秩算子都是紧算子, 因此我们有下述结论.

\begin{corollary}
\label{cor:6.2}
设 $(T_n)$ 为有限秩算子序列, $T\in\LL(E,F)$ 满足
$\norm{T_n-T}_{\LL(E,F)}\to0$.则 $T\in\KK(E,F)$.
\end{corollary}

\begin{starremark}
\label{rem:6.1}
著名的「逼近问题」(Banach, Grothendieck) 讨论的是推论  \ref{cor:6.2} 的逆:
给定紧算子 $T$, 是否总存在有限秩算子序列 $(T_n)$, 使得
$\norm{T_n-T}_{\LL(E,F)}\to0$? 这个问题长期悬而未决, 直到 1972 年
P.~Enflo \bref{Enflo}{1} 找到了反例.原来的构造相当复杂, 其后又找到了更简单的
例子, 例如取 $F$ 为 $\ell^p$ 的某个闭子空间 (对任意 $1<p<\infty$, $p\ne2$).
有兴趣的读者可在 J.~Lindenstrauss--L.~Tzafriri \bref{LindenstraussTzafriri}{2}
中找到关于逼近问题的详细讨论.注意, 在某些特殊情形下逼近问题的答案是肯定的
——例如当 $F$ 是 Hilbert 空间时.事实上, 令 $K=T(B_E)$.给定 $\varepsilon>0$,
存在 $K$ 的半径 $\varepsilon$ 的有限覆盖, 即
$K\subset\bigcup_{i\in I}B(f_i,\varepsilon)$.设 $G$ 为诸 $f_i$ 张成的向量空间,
并令 $T_\varepsilon=P_GT$, 于是 $T_\varepsilon$ 是有限秩的.我们断言
$\norm{T_\varepsilon-T}_{\LL(E,F)}<2\varepsilon$.对每个 $x\in B_E$, 存在
$i_0\in I$, 使得
\begin{equation}
\norm{Tx-f_{i_0}}<\varepsilon.
\label{eq:6.1}
\end{equation}
从而
\begin{equation*}
\norm{P_GTx-P_Gf_{i_0}}<\varepsilon,
\end{equation*}
即
\begin{equation}
\norm{P_GTx-f_{i_0}}<\varepsilon.
\label{eq:6.2}
\end{equation}
合并 \eqref{eq:6.1} 与 \eqref{eq:6.2}, 得到
\begin{equation*}
\norm{P_GTx-Tx}<2\varepsilon\quad\forall x\in B_E,
\end{equation*}
即
\begin{equation*}
\norm{T_\varepsilon-T}_{\LL(E,F)}<2\varepsilon.
\end{equation*}
[更一般地可以看出: 若 $F$ 有 Schauder 基, 则对每个空间 $E$ 以及由 $E$ 到
$F$ 的每个紧算子, 逼近问题的答案都是肯定的.]

联系到逼近问题, 我们提一种在非线性分析中非常有用的技巧, 即用有限秩的
非线性映射来逼近一个连续映射 (线性或非线性).设 $X$ 为拓扑空间, $F$ 为
Banach 空间, $T:X\to F$ 为连续映射, 使得 $T(X)$ 在 $F$ 中有紧的闭包.
我们断言: 对每个 $\varepsilon>0$, 存在有限秩的连续映射 $T_\varepsilon:X\to F$,
使得
\begin{equation}
\norm{T_\varepsilon(x)-T(x)}<\varepsilon\quad\forall x\in X.
\label{eq:6.3}
\end{equation}
事实上, 由于 $K=T(X)$ 紧, 存在 $K$ 的有限覆盖, 即
$K\subset\bigcup_{i\in I}B(f_i,\varepsilon/2)$.令
\begin{equation*}
T_\varepsilon(x)=\frac{\sum_{i\in I}q_i(x)f_i}{\sum_{i\in I}q_i(x)},
\quad\text{其中}\ q_i(x)=\max\{\varepsilon-\norm{Tx-f_i},0\};
\end{equation*}
显然 $T_\varepsilon$ 满足 \eqref{eq:6.3}.

这种逼近方式非常有用, 例如可用它由 Brouwer 不动点定理推出 Schauder 不动点
定理 (见例如 K.~Deimling \bref{Deimling}{1}, A.~Granas--J.~Dugundji \bref{GranasDugundji}{1}, J.~Franklin \bref{Franklin}{1}, 以及习题  \ref{exer:6.26}).类似的构造与
Schauder 不动点定理相结合, 还被 Lomonosov 以出人意料的方式用来证明一大类
线性算子存在非平凡不变子空间 (见例如 C.~Pearcy \bref{Pearcy}{1}, N.~Akhiezer--I.~Glazman \bref{AkhiezerGlazman}{1}, A.~Granas--J.~Dugundji \bref{GranasDugundji}{1}, 以及问题
42).另一个基于 Schauder 不动点定理而有简单证明的线性结果
是 Krein--Rutman 定理 (见定理  \ref{thm:6.13} 与问题 41).
\end{starremark}

\begin{proposition}
\label{prop:6.3}
设 $E$, $F$, $G$ 为三个 Banach 空间.设 $T\in\LL(E,F)$,
$S\in\KK(F,G)$ [相应地, $T\in\KK(E,F)$, $S\in\LL(F,G)$].则
$S\circ T\in\KK(E,G)$.
\end{proposition}

证明是显然的.

\begin{theorem}[Schauder]
\label{thm:6.4}
若 $T\in\KK(E,F)$, 则 $T^\star\in\KK(F^\star,E^\star)$.反之亦然.
\end{theorem}

\begin{proof}
我们要证明 $T^\star(B_{F^\star})$ 在 $E^\star$ 中有紧的闭包.设 $(v_n)$ 为
$B_{F^\star}$ 中的序列.我们断言 $(T^\star(v_n))$ 有收敛子列.令
$K=T(B_E)$; 这是一个紧度量空间.考虑由
\begin{equation*}
H=\{\phi_n:x\in K\mapsto\dual{v_n}{x};\ n=1,2,\dots\}
\end{equation*}
定义的集 $H\subset C(K)$.Ascoli--Arzelà 定理 (定理  \ref{thm:4.25}) 的假设
成立.于是存在子列 (仍记作 $\phi_{n_k}$), 它在 $K$ 上一致收敛于某个连续函数
$\phi\in C(K)$.特别地, 我们有
\begin{equation*}
\sup_{u\in B_E}\abs{\dual{v_{n_k}}{Tu}-\phi(Tu)}\xrightarrow[k\to\infty]{}0.
\end{equation*}
从而
\begin{equation*}
\sup_{u\in B_E}\abs{\dual{v_{n_k}}{Tu}-\dual{v_{n_\ell}}{Tu}}
\xrightarrow[k,\ell\to\infty]{}0,
\end{equation*}
即 $\norm{T^\star v_{n_k}-T^\star v_{n_\ell}}_{E^\star}
\xrightarrow[k,\ell\to\infty]{}0$.于是 $T^\star v_{n_k}$ 在 $E^\star$ 中
收敛.

反之, 设 $T^\star\in\KK(F^\star,E^\star)$.由第一部分, 我们已经知道
$T^{\star\star}\in\KK(E^{\star\star},F^{\star\star})$.特别地,
$T^{\star\star}(B_E)$ 在 $F^{\star\star}$ 中有紧的闭包.但
$T(B_E)=T^{\star\star}(B_E)$, 且 $F$ 在 $F^{\star\star}$ 中闭.因此
$T(B_E)$ 在 $F$ 中有紧的闭包.
\end{proof}

\begin{remark}
\label{rem:6.2}
设 $E$ 与 $F$ 为两个 Banach 空间, $T\in\KK(E,F)$.若 $(u_n)$ 在 $E$ 中
弱收敛于 $u$, 则 $(Tu_n)$ 强收敛于 $Tu$.若 $E$ 自反, 逆命题也成立
(见习题  \ref{exer:6.7}).
\end{remark}

\section{Riesz--Fredholm 理论}
\label{sec:6.2}

我们首先给出一些有用的预备结果.

\begin{lemma}[Riesz 引理]
\label{lem:6.1}
设 $E$ 为赋范向量空间 (n.v.s.), $M\subset E$ 为闭线性子空间, 且 $M\ne E$.
则
\begin{equation*}
\forall\varepsilon>0\ \ \exists u\in E\ \text{使得}\ \norm{u}=1\
\text{且}\ \dist(u,M)\ge1-\varepsilon.
\end{equation*}
\end{lemma}

\begin{proof}
取 $v\in E$, $v\notin M$.由于 $M$ 闭, 故
\begin{equation*}
d=\dist(v,M)>0.
\end{equation*}
任取 $m_0\in M$, 使得
\begin{equation*}
d\le\norm{v-m_0}\le d/(1-\varepsilon).
\end{equation*}
则
\begin{equation*}
u=\frac{v-m_0}{\norm{v-m_0}}
\end{equation*}
满足所需性质.事实上, 对每个 $m\in M$, 我们有
\begin{equation*}
\norm{u-m}=\left\|\frac{v-m_0}{\norm{v-m_0}}-m\right\|
\ge\frac{d}{\norm{v-m_0}}\ge1-\varepsilon,
\end{equation*}
因为 $m_0+\norm{v-m_0}m\in M$.
\end{proof}

\begin{remark}
\label{rem:6.3}
若 $M$ 是有限维的 (或更一般地, 若 $M$ 自反), 则在引理  \ref{lem:6.1} 中可
取 $\varepsilon=0$.但一般说来这并不成立 (见习题  \ref{exer:1.17}).
\end{remark}

\begin{keytheorem}[Riesz]
\label{thm:6.5}
设 $E$ 为赋范向量空间, 且 $B_E$ 紧.则 $E$ 是有限维的.
\end{keytheorem}

\begin{proof}
反设 $E$ 是无穷维的.于是存在 $E$ 的有限维子空间序列 $(E_n)$, 使得
$E_{n-1}\subset E_n$ 且 $E_{n-1}\ne E_n$.由引理  \ref{lem:6.1}, 存在序列
$(u_n)$, $u_n\in E_n$, 使得 $\norm{u_n}=1$ 且
$\dist(u_n,E_{n-1})\ge1/2$.特别地, 当 $m<n$ 时
$\norm{u_n-u_m}\ge1/2$.于是 $(u_n)$ 没有收敛子列, 这与 $B_E$ 紧的假设
矛盾.
\end{proof}

\begin{keytheorem}[Fredholm 二择一定理]
\label{thm:6.6}
设 $T\in\KK(E)$.则
\begin{enumerate}[label=(\alph*), leftmargin=2.2em, itemsep=0.1em]
	\item $N(I-T)$ 是有限维的,
	\item $R(I-T)$ 闭, 且更确切地, $R(I-T)=N(I-T^\star)^\perp$,
	\item $N(I-T)=\{0\}\iff R(I-T)=E$,
	\item $\dim N(I-T)=\dim N(I-T^\star)$.
\end{enumerate}
\end{keytheorem}

\begin{remark}
\label{rem:6.4}
Fredholm 二择一定理讨论方程 $u-Tu=f$ 的可解性.它说:
\begin{itemize}[leftmargin=1.6em, itemsep=0.15em]
	\item 或者对每个 $f\in E$, 方程 $u-Tu=f$ 有唯一解,
	\item 或者齐次方程 $u-Tu=0$ 有 $n$ 个线性无关解, 而此时非齐次方程
	      $u-Tu=f$ 可解当且仅当 $f$ 满足 $n$ 个正交条件, 即
	      $f\in N(I-T^\star)^\perp$.
\end{itemize}
\end{remark}

\begin{remark}
\label{rem:6.5}
性质 (c) 在有限维空间中是我们熟悉的.若 $\dim E<\infty$, 则由 $E$ 到自身的
线性算子单射 (= 一一) 当且仅当它满射 (= 到上).然而在无穷维空间中, 一个有界
算子可以单射而不满射, 反之亦然, 例如 $\ell^2$ 中的右移位 (相应地, 左移位)
(见注记 \ref{rem:6.6}).因此, 断言 (c) 是形如 $I-T$ (其中 $T\in\KK(E)$)
的算子的一个值得注意的性质.
\end{remark}

\begin{proof}
\begin{enumerate}[label=(\alph*), leftmargin=2.2em, itemsep=0.3em]
	\item 令 $E_1=N(I-T)$.则 $B_{E_1}\subset T(B_E)$, 从而 $B_{E_1}$ 紧.
	      由定理  \ref{thm:6.5}, $E_1$ 必为有限维的.
	\item 设 $f_n=u_n-Tu_n\to f$.我们要证明 $f\in R(I-T)$.令
	      $d_n=\dist(u_n,N(I-T))$.由于 $N(I-T)$ 有限维, 存在
	      $v_n\in N(I-T)$, 使得 $d_n=\norm{u_n-v_n}$.我们有
	      \begin{equation}
	      f_n=(u_n-v_n)-T(u_n-v_n).
	      \label{eq:6.4}
	      \end{equation}
	      我们断言 $\norm{u_n-v_n}$ 保持有界.若不然, 则存在子列使得
	      $\norm{u_{n_k}}-v_{n_k}\to\infty$.令
	      $w_n=(u_n-v_n)/\norm{u_n-v_n}$.由 \eqref{eq:6.4} 看出
	      $w_{n_k}-Tw_{n_k}\to0$.再取子列 (仍简记作 $w_{n_k}$), 可设
	      $Tw_{n_k}\to z$.于是 $w_{n_k}\to z$ 且 $z\in N(I-T)$, 从而
	      $\dist(w_{n_k},N(I-T))\to0$.另一方面,
	      \begin{equation*}
	      \dist(w_n,N(I-T))
	      =\frac{\dist(u_n,N(I-T))}{\norm{u_n-v_n}}=1
	      \end{equation*}
	      (因为 $v_n\in N(I-T)$); 矛盾.

	      于是 $\norm{u_n-v_n}$ 保持有界, 而 $T$ 是紧算子, 故可抽出子列使得
	      $T(u_{n_k}-v_{n_k})$ 收敛于某个极限 $\ell$.由 \eqref{eq:6.4}
	      推出 $u_{n_k}-v_{n_k}\to f+\ell$.令 $g=f+\ell$, 则 $g-Tg=f$,
	      即 $f\in R(I-T)$.这就完成了算子 $(I-T)$ 有闭值域这一事实的证明.
	      因此可以应用定理  \ref{thm:2.19}, 推出
	      \begin{equation*}
	      R(I-T)=N(I-T^\star)^\perp,\qquad
	      R(I-T^\star)=N(I-T)^\perp.
	      \end{equation*}
	\item 我们先证蕴涵 $\Rightarrow$.反设
	      \begin{equation*}
	      E_1=R(I-T)\ne E.
	      \end{equation*}
	      则 $E_1$ 是 Banach 空间且 $T(E_1)\subset E_1$.于是
	      $T|_{E_1}\in\KK(E_1)$, 且 $E_2=(I-T)(E_1)$ 是 $E_1$ 的闭子空间.
	      又 $E_2\ne E_1$ (因为 $(I-T)$ 单射).令 $E_n=(I-T)^n(E)$, 得到
	      一列 (严格) 递减的闭子空间.用 Riesz 引理可构造序列 $(u_n)$, 使得
	      $u_n\in E_n$, $\norm{u_n}=1$ 且 $\dist(u_n,E_{n+1})\ge1/2$.我们有
	      \begin{equation*}
	      Tu_n-Tu_m=-(u_n-Tu_n)+(u_m-Tu_m)+(u_n-u_m).
	      \end{equation*}
	      注意若 $n>m$, 则 $E_{n+1}\subset E_n\subset E_{m+1}\subset E_m$,
	      因此
	      \begin{equation*}
	      -(u_n-Tu_n)+(u_m-Tu_m)+u_n\in E_{m+1}.
	      \end{equation*}
	      由此推出 $\norm{Tu_n-Tu_m}\ge\dist(u_m,E_{m+1})\ge1/2$.这不可能,
	      因为 $T$ 是紧算子.于是我们证明了 $R(I-T)=E$.

	      反之, 设 $R(I-T)=E$.由推论  \ref{cor:2.18}, 我们知道
	      $N(I-T^\star)=R(I-T)^\perp=\{0\}$.由于 $T^\star\in\KK(E^\star)$,
	      可对之应用上一步, 推出 $R(I-T^\star)=E^\star$.再次使用推论
	       \ref{cor:2.18}, 得到 $N(I-T)=R(I-T^\star)^\perp=\{0\}$.
	\item 令 $d=\dim N(I-T)$, $d^\star=\dim N(I-T^\star)$.我们先证
	      $d^\star\le d$.若不然, 设 $d<d^\star$.由于 $N(I-T)$ 有限维,
	      它在 $E$ 中有补空间 (见第~\ref{sec:2.4}~节例 1).于是存在由 $E$ 到 $N(I-T)$
	      的连续投影算子 $P$.另一方面, $R(I-T)=N(I-T^\star)^\perp$ 有有限
	      余维 $d^\star$ (见第~\ref{sec:2.4}~节例 2), 从而它在 $E$ 中有维数为 $d^\star$
	      的补空间, 记作 $F$.由于 $d<d^\star$, 存在线性映射
	      $\Lambda:N(I-T)\to F$, 它单射但不满射.令 $S=T+\Lambda\circ P$.则
	      $S\in\KK(E)$, 因为 $\Lambda\circ P$ 是有限秩的.

	      我们断言 $N(I-S)=\{0\}$.事实上, 若
	      \begin{equation*}
	      0=u-Su=(u-Tu)-(\Lambda\circ Pu),
	      \end{equation*}
	      则
	      \begin{equation*}
	      u-Tu=0\quad\text{且}\quad \Lambda\circ Pu=0,
	      \end{equation*}
	      即 $u\in N(I-T)$ 且 $\Lambda u=0$.因此 $u=0$.

	      对算子 $S$ 应用 (c), 得到 $R(I-S)=E$.这是荒谬的, 因为存在
	      $f\in F$ 使 $f\notin R(\Lambda)$, 于是方程 $u-Su=f$ 无解.

	      于是我们证明了 $d^\star\le d$.把这一事实应用于 $T^\star$, 得到
	      \begin{equation*}
	      \dim N(I-T^{\star\star})\le\dim N(I-T^\star)\le\dim N(I-T).
	      \end{equation*}
	      但 $N(I-T^{\star\star})\supset N(I-T)$, 因此 $d=d^\star$.
\end{enumerate}
\end{proof}

\section{紧算子的谱}
\label{sec:6.3}

下面是一些重要的定义.

\begin{definition}
设 $T\in\LL(E)$.\emph{\index{yujieji@预解集 (resolvent set)}
预解集} (记作 $\rho(T)$) 定义为
\begin{equation*}
\rho(T)=\{\lambda\in\R;\ (T-\lambda I)\ \text{是由}\ E\ \text{到}\ E\ \text{的双射}\}.
\end{equation*}
\emph{谱} (记作 $\sigma(T)$) 是预解集的补集, 即 $\sigma(T)=\R\setminus\rho(T)$.
若 $N(T-\lambda I)\ne\{0\}$, 则称实数 $\lambda$ 是 $T$ 的\emph{\index{tezhengzhi@特征值 (eigenvalue)}
特征值};
$N(T-\lambda I)$ 是对应的\emph{特征空间}.全部特征值之集记作 $EV(T)$.
\footnote{有些作者用 $\sigma_p(T)$ (点谱) 代替 $EV(T)$.}
\end{definition}

记住下面一点是有用的: 若 $\lambda\in\rho(T)$, 则
$(T-\lambda I)^{-1}\in\LL(E)$ (见推论  \ref{cor:2.7}).

\begin{remark}
\label{rem:6.6}
显然 $EV(T)\subset\sigma(T)$.一般说来, 这个包含可以是严格的:
\footnote{当然, 若 $E$ 有限维, 则 $EV(T)=\sigma(T)$.}
可能存在某个 $\lambda$, 使得
\begin{equation*}
N(T-\lambda I)=\{0\}\quad\text{且}\quad R(T-\lambda I)\ne E
\end{equation*}
(这样的 $\lambda$ 属于谱但不是特征值).例如考虑 $E=\ell^2$ 中的右移位,
即 $Tu=(0,u_1,u_2,\dots)$, 其中 $u=(u_1,u_2,u_3,\dots)$.则
$0\in\sigma(T)$, 而 $0\notin EV(T)$.事实上, 此时 $EV(T)=\emptyset$, 而
$\sigma(T)=[-1,+1]$ (见习题  \ref{exer:6.18}).当然, 在有限维或无穷维空间中都
可能出现 $EV(T)=\sigma(T)=\emptyset$; 例如考虑 $\R^2$ 中旋转 $\pi/2$ 的
变换, 或在 $\ell^2$ 中考虑算子 $Tu=(-u_2,u_1,-u_4,u_3,\dots)$.若在 $\CC$
上的向量空间中讨论 (见第~\ref{sec:11.4}~节), 情况就完全不同了; 在 $\CC$ 上的空间中,
特征值与谱的研究要有趣得多.众所周知, 在 $\CC$ 上的有限维空间中,
$EV(T)=\sigma(T)\ne\emptyset$ (它们是特征多项式的根).在 $\CC$ 上的无穷维
空间中, 一个非平凡的结果断言 $\sigma(T)$ 总是非空的 (见第~\ref{sec:11.4}~节).
然而, $EV(T)=\emptyset$ 是可能发生的 (例如取 $E=\ell^2$ 中的右移位).
\end{remark}

\begin{proposition}
\label{prop:6.7}
有界算子 $T$ 的谱 $\sigma(T)$ 紧, 且
$\sigma(T)\subset[-\norm{T},+\norm{T}]$.
\end{proposition}

\begin{proof}
设 $\lambda\in\R$ 满足 $\abs{\lambda}>\norm{T}$.我们将证明 $T-\lambda I$ 是
双射, 从而 $\sigma(T)\subset[-\norm{T},+\norm{T}]$.给定 $f\in E$, 方程
$Tu-\lambda u=f$ 有唯一解, 因为它可写成 $u=\lambda^{-1}(Tu-f)$, 而压缩映射
原理 (定理  \ref{thm:5.7}) 适用.

下面证明 $\rho(T)$ 是开集.设 $\lambda_0\in\rho(T)$.给定 $\lambda\in\R$
(接近 $\lambda_0$) 与 $f\in E$, 我们试图求解
\begin{equation}
Tu-\lambda u=f.
\label{eq:6.5}
\end{equation}
方程 \eqref{eq:6.5} 可写成
\begin{equation*}
Tu-\lambda_0u=f+(\lambda-\lambda_0)u,
\end{equation*}
即
\begin{equation}
u=(T-\lambda_0I)^{-1}\bigl[f+(\lambda-\lambda_0)u\bigr].
\label{eq:6.6}
\end{equation}
再次应用压缩映射原理, 可见当
$\abs{\lambda-\lambda_0}\,\norm{(T-\lambda_0I)^{-1}}<1$ 时 \eqref{eq:6.6}
有解.
\end{proof}

\begin{keytheorem}
\label{thm:6.8}
设 $T\in\KK(E)$ 且 $\dim E=\infty$, 则我们有:
\begin{enumerate}[label=(\alph*), leftmargin=2.2em, itemsep=0.1em]
	\item $0\in\sigma(T)$,
	\item $\sigma(T)\setminus\{0\}=EV(T)\setminus\{0\}$,
	\item 下述情形之一成立:
	      \begin{itemize}[leftmargin=1.6em, itemsep=0.15em]
	      	\item $\sigma(T)=\{0\}$,
	      	\item $\sigma(T)\setminus\{0\}$ 是有限集,
	      	\item $\sigma(T)\setminus\{0\}$ 是收敛于 $0$ 的序列.
	      \end{itemize}
\end{enumerate}
\end{keytheorem}

\begin{proof}
\begin{enumerate}[label=(\alph*), leftmargin=2.2em, itemsep=0.3em]
	\item 若不然, 设 $0\notin\sigma(T)$.则 $T$ 是双射, 而
	      $I=T\circ T^{-1}$ 是紧的.于是 $B_E$ 紧, 由定理  \ref{thm:6.5} 得
	      $\dim E<\infty$; 矛盾.
	\item 设 $\lambda\in\sigma(T)$, $\lambda\ne0$.我们要证明 $\lambda$ 是
	      特征值.若不然, 设 $N(T-\lambda I)=\{0\}$.则由定理  \ref{thm:6.6}(c),
	      我们知道 $R(T-\lambda I)=E$, 因此 $\lambda\in\rho(T)$; 矛盾.
\end{enumerate}

为证明断言 (c), 我们将使用下述引理.

\begin{lemma}
\label{lem:6.2}
设 $T\in\KK(E)$, $(\lambda_n)_{n\ge1}$ 为互不相同的实数序列, 使得
\begin{equation*}
\lambda_n\to\lambda\quad\text{且}\quad
\lambda_n\in\sigma(T)\setminus\{0\}\ \ \forall n.
\end{equation*}
则 $\lambda=0$.换言之, $\sigma(T)\setminus\{0\}$ 的所有点都是孤立点.
\end{lemma}

\begin{proof}
我们知道 $\lambda_n\in EV(T)$; 取 $e_n\ne0$ 使 $(T-\lambda_nI)e_n=0$.令
$E_n$ 为由 $\{e_1,e_2,\dots,e_n\}$ 张成的空间.我们断言对一切 $n$ 都有
$E_n\subset E_{n+1}$, $E_n\ne E_{n+1}$.只须验证对一切 $n$, 向量
$e_1,e_2,\dots,e_n$ 线性无关.对 $n$ 归纳证明.设此结论对 $n$ 成立, 并设
$e_{n+1}=\sum_{i=1}^n\alpha_ie_i$.则
\begin{equation*}
Te_{n+1}=\sum_{i=1}^n\alpha_i\lambda_ie_i
=\sum_{i=1}^n\alpha_i\lambda_{n+1}e_i.
\end{equation*}
由此推出 $\alpha_i(\lambda_i-\lambda_{n+1})=0$, $i=1,2,\dots,n$, 从而
$\alpha_i=0$, $i=1,2,\dots,n$; 矛盾.于是我们证明了
$E_n\subset E_{n+1}$, $E_n\ne E_{n+1}$ 对一切 $n$ 成立.

应用 Riesz 引理 (引理  \ref{lem:6.1}), 可构造序列 $(u_n)_{n\ge1}$, 使得
$u_n\in E_n$, $\norm{u_n}=1$ 且 $\dist(u_n,E_{n-1})\ge1/2$ 对一切 $n\ge2$
成立.当 $2\le m<n$ 时, 有
\begin{equation*}
E_{m-1}\subset E_m\subset E_{n-1}\subset E_n.
\end{equation*}
另一方面, 显然 $(T-\lambda_nI)E_n\subset E_{n-1}$.于是
\begin{equation*}
\left\|\frac{Tu_n}{\lambda_n}-\frac{Tu_m}{\lambda_m}\right\|
=\left\|\frac{(Tu_n-\lambda_nu_n)}{\lambda_n}
-\frac{(Tu_m-\lambda_mu_m)}{\lambda_m}+u_n-u_m\right\|
\ge\dist(u_n,E_{n-1})\ge1/2.
\end{equation*}
若 $\lambda_n\to\lambda$ 且 $\lambda\ne0$, 则矛盾, 因为 $(Tu_n)$ 有收敛子列.
\end{proof}

定理  \ref{thm:6.8} 的证明 (续).对每个整数 $n\ge1$, 集
\begin{equation*}
\sigma(T)\cap\{\lambda\in\R;\ \abs{\lambda}\ge1/n\}
\end{equation*}
或空或有限 (若它含无穷多个互不相同的点, 则存在子列收敛于某个满足
$\abs{\lambda}\ge1/n$ 的 $\lambda$——因为 $\sigma(T)$ 紧——这与引理
 \ref{lem:6.2} 矛盾).因此, 若 $\sigma(T)\setminus\{0\}$ 含无穷多个互不相同的
点, 我们可把它们排成一个趋于 $0$ 的序列.
\end{proof}

\begin{remark}
\label{rem:6.7}
给定任何收敛于 $0$ 的序列 $(\alpha_n)$, 存在紧算子 $T$ 使得
$\sigma(T)=(\alpha_n)\cup\{0\}$.在 $\ell^2$ 中, 只须考虑由
$Tu=(\alpha_1u_1,\alpha_2u_2,\dots,\alpha_nu_n,\dots)$ (其中
$u=(u_1,u_2,\dots,u_n,\dots)$) 定义的乘法算子 $T$.注意 $T$ 是紧的,
因为 $T$ 是有限秩算子的极限.更确切地说, 令
$T_nu=(\alpha_1u_1,\alpha_2u_2,\dots,\alpha_nu_n,0,0,\dots)$; 则
$\norm{T_n-T}\to0$.在这个例子中我们还看到, $0$ 可以属于也可以不属于
$EV(T)$.另一方面, 若 $0\in EV(T)$, 则对应的特征空间, 即 $N(T)$, 可以是
有限维或无穷维的.
\end{remark}

\section{自伴紧算子的谱分解}
\label{sec:6.4}

下文假设 $E=H$ 是 Hilbert 空间, $T\in\LL(H)$.把 $H^\star$ 与 $H$ 等同,
我们可把 $T^\star$ 看作由 $H$ 到自身的算子.

\begin{definition}
若 $T^\star=T$, 即
\begin{equation*}
(Tu,v)=(u,Tv)\quad\forall u,v\in H,
\end{equation*}
则称有界算子 $T\in\LL(H)$ 是\emph{自伴}的.
\end{definition}

\begin{proposition}
\label{prop:6.9}
设 $T\in\LL(H)$ 为自伴算子.令
\begin{equation*}
m=\inf_{\substack{u\in H\\ \abs{u}=1}}(Tu,u)
\quad\text{与}\quad
M=\sup_{\substack{u\in H\\ \abs{u}=1}}(Tu,u).
\end{equation*}
则 $\sigma(T)\subset[m,M]$, $m\in\sigma(T)$, $M\in\sigma(T)$.此外,
$\norm{T}=\max\{\abs{m},\abs{M}\}$.
\end{proposition}

\begin{proof}
设 $\lambda>M$; 我们将证明 $\lambda\in\rho(T)$.我们有
\begin{equation*}
(Tu,u)\le M\abs{u}^2\quad\forall u\in H,
\end{equation*}
从而
\begin{equation*}
(\lambda u-Tu,u)\ge(\lambda-M)\abs{u}^2=\alpha\abs{u}^2
\quad\forall u\in H,\ \text{其中}\ \alpha>0.
\end{equation*}
应用 Lax--Milgram 定理 (推论  \ref{cor:5.8}), 推出 $\lambda I-T$ 是双射,
因此 $\lambda\in\rho(T)$.类似地, 任何 $\lambda<m$ 都属于 $\rho(T)$, 因此
$\sigma(T)\subset[m,M]$.

下面证明 $M\in\sigma(T)$ (证明 $m\in\sigma(T)$ 类似).双线性形式
$a(u,v)=(Mu-Tu,v)$ 对称且满足
\begin{equation*}
a(v,v)\ge0\quad\forall v\in H.
\end{equation*}
因此它满足 Cauchy--Schwarz 不等式
\begin{equation*}
\abs{a(u,v)}\le a(u,u)^{1/2}a(v,v)^{1/2}\quad\forall u,v\in H,
\end{equation*}
即
\begin{equation*}
\abs{(Mu-Tu,v)}\le(Mu-Tu,u)^{1/2}(Mv-Tv,v)^{1/2}
\quad\forall u,v\in H.
\end{equation*}
由此推出
\begin{equation}
\abs{Mu-Tu}\le C(Mu-Tu,u)^{1/2}\quad\forall u\in H.
\label{eq:6.7}
\end{equation}
由 $M$ 的定义, 存在序列 $(u_n)$ 使得 $\abs{u_n}=1$ 且
$(Tu_n,u_n)\to M$.由 \eqref{eq:6.7} 推出 $\abs{Mu_n-Tu_n}\to0$, 从而
$M\in\sigma(T)$ (因为若 $M\in\rho(T)$, 则
$u_n=(MI-T)^{-1}(Mu_n-Tu_n)\to0$, 这不可能).

最后, 我们证明 $\norm{T}=\mu$, 其中 $\mu=\max\{\abs{m},\abs{M}\}$.对
$u,v\in H$ 写出
\begin{align*}
(T(u+v),u+v)&=(Tu,u)+(Tv,v)+2(Tu,v),\\
(T(u-v),u-v)&=(Tu,u)+(Tv,v)-2(Tu,v).
\end{align*}
于是
\begin{equation*}
4(Tu,v)=(T(u+v),u+v)-(T(u-v),u-v)
\le M\abs{u+v}^2-m\abs{u-v}^2,
\end{equation*}
从而
\begin{equation*}
4\abs{(Tu,v)}\le\mu\bigl(\abs{u+v}^2+\abs{u-v}^2\bigr)
=2\mu\bigl(\abs{u}^2+\abs{v}^2\bigr).
\end{equation*}
用 $\alpha>0$ 代替 $v$ 为 $\alpha v$, 得到
\begin{equation*}
4\abs{(Tu,v)}\le2\mu\left(\frac{\abs{u}^2}{\alpha}+\alpha\abs{v}^2\right).
\end{equation*}
下一步对 $\alpha$ 极小化右端, 即取 $\alpha=\abs{u}/\abs{v}$, 然后得到
\begin{equation*}
\abs{(Tu,v)}\le\mu\abs{u}\,\abs{v}\quad\forall u,v,
\end{equation*}
从而 $\norm{T}\le\mu$.另一方面, 显然 $\abs{(Tu,u)}\le\norm{T}\abs{u}^2$,
故 $\abs{m}\le\norm{T}$ 且 $\abs{M}\le\norm{T}$, 从而 $\mu\le\norm{T}$.
\end{proof}

\begin{corollary}
\label{cor:6.10}
设 $T\in\LL(H)$ 为自伴算子, 且 $\sigma(T)=\{0\}$.则 $T=0$.
\end{corollary}

我们最后的结论是一个基本结果.它断言每个紧自伴算子都可以在某个适当的基下
对角化.

\begin{keytheorem}
\label{thm:6.11}
设 $H$ 为可分 Hilbert 空间, $T$ 为紧自伴算子.则存在由 $T$ 的特征向量构成的
Hilbert 基.
\end{keytheorem}

\begin{proof}
设 $(\lambda_n)_{n\ge1}$ 为 $T$ 的全部 (互不相同的) 非零特征值.令
\begin{equation*}
\lambda_0=0,\qquad E_0=N(T),\qquad E_n=N(T-\lambda_nI).
\end{equation*}
回忆
\begin{equation*}
0\le\dim E_0\le\infty\quad\text{且}\quad0<\dim E_n<\infty.
\end{equation*}
我们断言 $H$ 是诸 $E_n$ ($n=0,1,2,\dots$) 的 Hilbert 和 (按第~\ref{sec:5.4}~节的
意义):

(i) 空间 $(E_n)_{n\ge0}$ 两两正交.

事实上, 若 $u\in E_m$, $v\in E_n$ 且 $m\ne n$, 则
\begin{equation*}
Tu=\lambda_mu\quad\text{且}\quad Tv=\lambda_nv,
\end{equation*}
从而
\begin{equation*}
(Tu,v)=\lambda_m(u,v)=(u,Tv)=\lambda_n(u,v).
\end{equation*}
因此 $(u,v)=0$.

(ii) 设 $F$ 为诸空间 $(E_n)_{n\ge0}$ 张成的向量空间.我们将证明 $F$ 在 $H$
中稠密.

显然 $T(F)\subset F$.由此推出 $T(F^\perp)\subset F^\perp$; 事实上, 给定
$u\in F^\perp$, 我们有
\begin{equation*}
(Tu,v)=(u,Tv)=0\quad\forall v\in F,
\end{equation*}
故 $Tu\in F^\perp$.$T$ 在 $F^\perp$ 上的限制记作 $T_0$.它是 $F^\perp$ 上的
紧自伴算子.我们断言 $\sigma(T_0)=\{0\}$.若不然, 设某个 $\lambda\ne0$ 属于
$\sigma(T_0)$.由于 $\lambda\in EV(T_0)$, 存在 $u\in F^\perp$, $u\ne0$, 使得
$T_0u=\lambda u$.于是 $\lambda$ 是 $T$ 的特征值之一, 设为
$\lambda=\lambda_n$ ($n\ge1$).从而 $u\in E_n\subset F$.由于
$u\in F^\perp\cap F$, 推出 $u=0$; 矛盾.应用推论  \ref{cor:6.10}, 推出
$T_0=0$, 即 $T$ 在 $F^\perp$ 上为零.由此推出 $F^\perp\subset N(T)$.另一方面,
$N(T)\subset F$, 因而 $F^\perp\subset F$.这意味着 $F^\perp=\{0\}$, 故 $F$
在 $H$ 中稠密.

最后, 在每个子空间 $(E_n)_{n\ge0}$ 中选取一个 Hilbert 基 ($E_0$ 的这种基的
存在性由定理  \ref{thm:5.11} 得出; 对其余的 $E_n$, $n\ge1$, 这是显然的,
因为它们有限维).这些基的并显然是 $H$ 的由 $T$ 的特征向量构成的 Hilbert 基.
\end{proof}

\begin{remark}
\label{rem:6.8}
设 $T$ 为紧自伴算子.由前面的分析, 我们可以把任何 $u\in H$ 写成
\begin{equation*}
u=\sum_{n=0}^\infty u_n\quad\text{其中}\ u_n\in E_n.
\end{equation*}
则 $Tu=\sum_{n=1}^\infty\lambda_nu_n$.给定整数 $k\ge1$, 令
\begin{equation*}
T_ku=\sum_{n=1}^k\lambda_nu_n.
\end{equation*}
显然 $T_k$ 是有限秩算子, 且
\begin{equation*}
\norm{T_k-T}\le\sup_{n\ge k+1}\abs{\lambda_n}\to0
\quad\text{当}\ k\to\infty.
\end{equation*}
回忆: 事实上在 Hilbert 空间中, 每个紧算子——不必自伴——都是有限秩算子序列的
极限 (见注记 \ref{rem:6.1}).
\end{remark}

\section*{第 6 章评注}
\addcontentsline{toc}{section}{第 6 章评注}
\markright{第 6 章评注}

\cmtitem{$\star$1. Fredholm 算子.}
定理  \ref{thm:6.6} 是通向 Fredholm 算子理论的第一步.给定两个 Banach 空间
$E$ 与 $F$, 若 $A\in\LL(E,F)$ 满足
\begin{enumerate}[label=(\roman*), leftmargin=2.6em, itemsep=0.1em]
	\item $N(A)$ 是有限维的,
	\item $R(A)$ 闭且余维有限,\footnote{设 $A\in\LL(E,F)$ 满足 $N(A)$ 有限维
	      且 $R(A)$ 余维有限 (即存在有限维空间 $G\subset F$, 使得
	      $R(A)+G=F$).则可推出 $R(A)$ 闭 (见习题  \ref{exer:2.27}).}
\end{enumerate}
则称 $A$ 是\emph{Fredholm 算子} (或 \emph{Noether 算子})——记作
$A\in\Phi(E,F)$.$A$ 的\emph{指标}定义为
\begin{equation*}
\ind A=\dim N(A)-\codim R(A).
\end{equation*}
例如, $A=I-T$ (其中 $T\in\KK(E)$) 是指标为零的 Fredholm 算子; 这由定理
 \ref{thm:6.6} 得出.

Fredholm 算子的主要性质如下:

\begin{enumerate}[label=(\alph*), leftmargin=2.2em, itemsep=0.3em]
	\item Fredholm 算子类 $\Phi(E,F)$ 是 $\LL(E,F)$ 的开子集, 且映射
	      $A\mapsto\ind A$ 连续; 因此它在 $\Phi(E,F)$ 的每个连通分支上为常数.
	\item 每个算子 $A\in\Phi(E,F)$ 模有限秩算子可逆, 即存在算子
	      $B\in\LL(F,E)$, 使得
	      \begin{equation*}
	      (A\circ B-I_F)\quad\text{与}\quad(B\circ A-I_E)
	      \end{equation*}
	      都是有限秩算子.

	      反之, 设 $A\in\LL(E,F)$, 并设存在 $B\in\LL(F,E)$, 使得
	      \begin{equation*}
	      A\circ B-I_F\in\KK(F)\quad\text{且}\quad
	      B\circ A-I_E\in\KK(E).
	      \end{equation*}
	      则 $A\in\Phi(E,F)$.
	\item 若 $A\in\Phi(E,F)$ 且 $T\in\KK(E,F)$, 则 $A+T\in\Phi(E,F)$ 且
	      $\ind(A+T)=\ind A$.
	\item 若 $A\in\Phi(E,F)$ 且 $B\in\Phi(F,G)$, 则 $B\circ A\in\Phi(E,G)$ 且
	      $\ind(B\circ A)=\ind(A)+\ind(B)$.
\end{enumerate}

关于这个问题, 见例如 T.~Kato \bref{Kato}{1}, M.~Schechter \bref{Schechter}{1}, S.~Lang \bref{Lang}{1}, A.~E.~Taylor--D.~C.~Lay \bref{TaylorLay}{1}, P.~Lax \bref{Lax}{1}, L.~Hörmander \bref{Hormander}{2} (第 3 卷), 以及问题
38.

\cmtitem{$\star$2. Hilbert--Schmidt 算子.}
设 $H$ 为可分 Hilbert 空间.若有 $H$ 中的 Hilbert 基 $(e_n)$, 使得
$\norm{T}_{HS}^2=\sum\abs{Te_n}^2<\infty$, 则称有界算子 $T\in\LL(H)$ 为
\emph{Hilbert--Schmidt 算子}.可以证明这个定义与基的选取无关, 且
$\norm{\cdot}_{HS}$ 是一个范数.每个 Hilbert--Schmidt 算子都是紧的.
Hilbert--Schmidt 算子起着重要作用, 尤其是由于下述结果:

\begin{theorem}
\label{thm:6.12}
设 $H=L^2(\Omega)$, $K(x,y)\in L^2(\Omega\times\Omega)$.则算子
\begin{equation*}
u\mapsto(Ku)(x)=\int_\Omega K(x,y)u(y)\,dy
\end{equation*}
是 Hilbert--Schmidt 算子.

反之, $L^2(\Omega)$ 上的每个 Hilbert--Schmidt 算子都具有上述形式, 其中
$K(x,y)\in L^2(\Omega\times\Omega)$ 由算子唯一确定.
\end{theorem}

关于这个问题, 见例如 A.~Balakrishnan \bref{Balakrishnan}{1}, N.~Dunford--J.~T.~Schwartz \bref{DunfordSchwartz}{1} 第 2 卷, 以及问题 40.

\cmtitem{3. 特征值的重数.}
设 $T\in\KK(E)$, $\lambda\in\sigma(T)\setminus\{0\}$.可以证明序列
$N((T-\lambda I)^k)$, $k=1,2,\dots$, 严格递增到某个有限的 $p$ 之后保持
不变 (见例如 A.~E.~Taylor--D.~C.~Lay \bref{TaylorLay}{1}, E.~Kreyszig \bref{Kreyszig}{1}, 以及问题 36).这个整数 $p$ 称为
$(T-\lambda I)$ 的\emph{升数}.$N(T-\lambda I)$ 的维数被某些作者称为
$\lambda$ 的\emph{几何重数}, $N((T-\lambda I)^p)$ 的维数称为 $\lambda$ 的
\emph{代数重数}; 若 $E$ 是 Hilbert 空间且 $T$ 自伴, 二者相等 (见问题 36).

\cmtitem{4. 谱分析.}
设 $H$ 为 Hilbert 空间, $T\in\LL(H)$ 为自伴算子 (可以不是紧的).
有一种称为 $T$ 的\emph{谱族}的构造, 它推广了第~\ref{sec:6.4}~节的谱分解.它特别允许
定义\emph{泛函演算}, 即对任何连续函数 $f$ 赋予量 $f(T)$ 意义.它还可以推广到
无界算子与非自伴算子, 只要假设 $T$ 是\emph{正规}的, 即 $TT^\star=T^\star T$.
谱分析是一个广阔的领域, 尤其是在 $\CC$ 上的 Banach 空间中 (见第~\ref{sec:11.4}~节),
有许多应用与分支.初等的介绍见例如 W.~Rudin \bref{Rudin}{1}, E.~Kreyszig \bref{Kreyszig}{1}, A.~Friedman \bref{Friedman}{3}, K.~Yosida \bref{Yosida}{1}.更完整的叙述见例如 M.~Reed--B.~Simon
\bref{ReedSimon}{1}, T.~Kato \bref{Kato}{1}, R.~Dautray--J.-L.~Lions \bref{DautrayLions}{1} 第 VIII 与 IX 章, N.~Dunford--J.~T.~Schwartz \bref{DunfordSchwartz}{1} 第 2 卷, N.~Akhiezer--I.~Glazman \bref{AkhiezerGlazman}{1}, A.~E.~Taylor--D.~C.~Lay \bref{TaylorLay}{1}, J.~Weidmann \bref{Weidmann}{1}, J.~B.~Conway \bref{Conway}{1}, P.~Lax \bref{Lax}{1}, 以及
M.~Schechter \bref{Schechter}{2}.

\cmtnum{5. 极小极大原理.} Courant--Fischer 的极小极大公式提供了计算特征值的
一种非常有用的方法; 见例如 R.~Courant--D.~Hilbert \bref{CourantHilbert}{1}, P.~Lax \bref{Lax}{1}, 以及问题 37.H.~Weinberger \bref{Weinberger}{2} 的
专著包含这个专题的大量发展.

\cmtitem{6. Krein--Rutman 定理.}
下述结果在研究二阶椭圆算子的谱性质时有重要应用 (见第~
\ref{chap:9}~章).

\begin{startheorem}[Krein--Rutman]
\label{thm:6.13}
设 $E$ 为 Banach 空间, $P$ 为顶点在 $0$ 的凸锥, 即
$\lambda x+\mu y\in P\ \forall\lambda\ge0$, $\forall\mu\ge0$,
$\forall x\in P$, $\forall y\in P$.假设 $P$ 闭, $\Int P\ne\emptyset$,
且 $P\ne E$.设 $T\in\KK(E)$ 满足 $T(P\setminus\{0\})\subset\Int P$.则存在
$x_0\in\Int P$ 与 $\lambda_0>0$, 使得 $Tx_0=\lambda_0x_0$; 此外, $\lambda_0$
是唯一对应于 $P$ 中特征向量的特征值, 即若 $Tx=\lambda x$, $x\in P$,
$x\ne0$, 则 $\lambda=\lambda_0$ 且 $x=mx_0$ (其中 $m>0$).最后,
\begin{equation*}
\lambda_0=\max\{\abs{\lambda};\ \lambda\in\sigma(T)\},
\end{equation*}
且 $\lambda$ 的重数 (几何重数与代数重数) 等于一.
\end{startheorem}

问题 41 中给出的证明属于 P.~Rabinowitz \bref{Rabinowitz}{2}.上述
Krein--Rutman 定理的变体可在例如 H.~Schaefer \bref{Schaefer}{1}, R.~Nussbaum \bref{Nussbaum}{1}, F.~F.~Bonsall \bref{Bonsall}{1} 与
J.~F.~Toland \bref{Toland}{4} 中找到.

\section*{第 6 章习题}
\addcontentsline{toc}{section}{第 6 章习题}
\markright{第 6 章习题}

\begin{exercise}
\label{exer:6.1}
设 $E=\ell^p$, $1\le p\le\infty$ (见第~\ref{sec:11.3}~节).设 $(\lambda_n)$ 为 $\R$
中的有界序列, 考虑由
\begin{equation*}
Tx=(\lambda_1x_1,\lambda_2x_2,\dots,\lambda_nx_n,\dots)
\end{equation*}
定义的算子 $T\in\LL(E)$, 其中 $x=(x_1,x_2,\dots,x_n,\dots)$.证明 $T$ 是由
$E$ 到 $E$ 的紧算子当且仅当 $\lambda_n\to0$.
\end{exercise}

\begin{exercise}
\label{exer:6.2}
设 $E$ 与 $F$ 为两个 Banach 空间, $T\in\LL(E,F)$.
\begin{enumerate}[label=\arabic*., leftmargin=2.2em, itemsep=0.3em]
	\item 假设 $E$ 自反.证明 $T(B_E)$ 是 (强) 闭的.
	\item 假设 $E$ 自反且 $T\in\KK(E,F)$.证明 $T(B_E)$ 紧.
	\item 设 $E=F=C([0,1])$, $Tu(t)=\int_0^tu(s)\,ds$.验证 $T\in\KK(E)$.
	      证明 $T(B_E)$ 不闭.
\end{enumerate}
\end{exercise}

\begin{exercise}
\label{exer:6.3}
设 $E$ 与 $F$ 为两个 Banach 空间, $T\in\KK(E,F)$.假设 $\dim E=\infty$.
证明存在 $E$ 中的序列 $(u_n)$, 使得 $\norm{u_n}_E=1$ 且
$\norm{Tu_n}_F\to0$.
[提示: 用反证法.]
\end{exercise}

\begin{exercise}
\label{exer:6.4}
设 $1\le p<\infty$.验证 $\ell^p\subset c_0$ 且嵌入连续 (关于 $\ell^p$ 与
$c_0$ 的定义, 见第~\ref{sec:11.3}~节).这个嵌入是紧的吗?
[提示: 使用 $\ell^p$ 的典范基 $(e_n)$.]
\end{exercise}

\begin{exercise}
\label{exer:6.5}
设 $(\lambda_n)$ 为正数序列, 满足 $\lim_{n\to\infty}\lambda_n=+\infty$.
设 $V$ 为满足
\begin{equation*}
\sum_{n=1}^\infty\lambda_n\abs{u_n}^2<\infty
\end{equation*}
的序列 $(u_n)_{n\ge1}$ 构成的空间.$V$ 赋以标量积
\begin{equation*}
((u,v))=\sum_{n=1}^\infty\lambda_nu_nv_n.
\end{equation*}
证明 $V$ 是 Hilbert 空间, 且 $V\subset\ell^2$ 的嵌入是紧的.
\end{exercise}

\begin{exercise}
\label{exer:6.6}
设 $1\le q\le p\le\infty$.证明由 $L^p(0,1)$ 到 $L^q(0,1)$ 的典范嵌入
连续但不紧.
[提示: 使用 Rademacher 函数; 见习题  \ref{exer:4.18}.]
\end{exercise}

\begin{exercise}
\label{exer:6.7}
设 $E$ 与 $F$ 为两个 Banach 空间, $T\in\LL(E,F)$.考虑下述性质:
\begin{equation*}
\text{(P)}\qquad
\begin{cases}
\text{对 }E\text{ 中每个弱收敛序列 }(u_n),\\
u_n\weakto u\ \text{蕴含}\ Tu_n\to Tu\ \text{在 }F\text{ 中强收敛}.
\end{cases}
\end{equation*}
\begin{equation*}
\text{(Q)}\qquad
\begin{cases}
T\text{ 是由 }E\text{ (赋以弱拓扑 }\sigma(E,E^\star)\text{) 到 }F\\
\text{(赋以强拓扑) 的连续算子}.
\end{cases}
\end{equation*}
\begin{enumerate}[label=\arabic*., leftmargin=2.2em, itemsep=0.3em]
	\item 证明
	      \begin{equation*}
	      \text{(Q)}\iff T\ \text{是有限秩算子}.
	      \end{equation*}
	\item 证明 $T\in\KK(E,F)\Rightarrow\text{(P)}$.
	\item 假设 $E=\ell^1$ 或 $F=\ell^1$.证明每个算子 $T\in\LL(E,F)$ 都满足
	      (P).[提示: 使用问题 8 的一个结果.]
\end{enumerate}

下文假设 $E$ 自反.

\begin{enumerate}[label=\arabic*., leftmargin=2.2em, start=4, itemsep=0.3em]
	\item 证明 $T\in\KK(E,F)\iff\text{(P)}$.
	\item 由此推出每个算子 $T\in\LL(E,\ell^1)$ 都是紧的.
	\item 证明每个算子 $T\in\LL(c_0,E)$ 都是紧的.
	      [提示: 考虑伴随算子 $T^\star$.]
\end{enumerate}
\end{exercise}

\begin{exercise}
\label{exer:6.8}
设 $E$ 与 $F$ 为两个 Banach 空间, $T\in\KK(E,F)$.假设 $R(T)$ 闭.
\begin{enumerate}[label=\arabic*., leftmargin=2.2em, itemsep=0.3em]
	\item 证明 $T$ 是有限秩算子.
	      [提示: 使用开映射定理, 即定理  \ref{thm:2.6}.]
	\item 此外假设 $\dim N(T)<\infty$.证明 $\dim E<\infty$.
\end{enumerate}
\end{exercise}

\begin{exercise}
\label{exer:6.9}
设 $E$ 与 $F$ 为两个 Banach 空间, $T\in\LL(E,F)$.
\begin{enumerate}[label=\arabic*., leftmargin=2.2em, itemsep=0.3em]
	\item 证明下述三个性质等价:\footnote{投影算子是指满足 $P^2=P$ 的
	      算子 $P$.}
	      \begin{equation*}
	      \text{(A)}\qquad
	      \begin{cases}
	      \dim N(T)<\infty\ \text{且}\ R(T)\ \text{闭}.
	      \end{cases}
	      \end{equation*}
	      \begin{equation*}
	      \text{(B)}\qquad
	      \begin{cases}
	      \text{存在有限秩投影算子 }P\in\LL(E)\text{ 与常数 }C,\\
	      \text{使得}\ \norm{u}_E\le C(\norm{Tu}_F+\norm{Pu}_E)
	      \quad\forall u\in E.
	      \end{cases}
	      \end{equation*}
	      \begin{equation*}
	      \text{(C)}\qquad
	      \begin{cases}
	      \text{存在 Banach 空间 }G\text{, 算子 }Q\in\KK(E,G)\text{ 与常数 }C,\\
	      \text{使得}\ \norm{u}_E\le C(\norm{Tu}_F+\norm{Qu}_G)
	      \quad\forall u\in E.
	      \end{cases}
	      \end{equation*}
	      [提示: 当 $\dim N(T)<\infty$ 时, 考虑 $N(T)$ 的补空间;
	      见第~\ref{sec:2.4}~节.] 与习题  \ref{exer:2.12} 作比较.
	\item 假设 $T$ 满足 (A).证明对每个 $S\in\KK(E,F)$, $(T+S)$ 也满足 (A).
	\item 证明满足 (A) 的全部算子 $T\in\LL(E,F)$ 之集在 $\LL(E,F)$ 中是开集.
	\item 设 $F_0$ 为 $F$ 的闭线性子空间, $S\in\KK(F_0,F)$.证明
	      $(I+S)(F_0)$ 是 $F$ 的闭子空间.
\end{enumerate}
\end{exercise}

\begin{exercise}
\label{exer:6.10}
设 $Q(t)=\sum_{k=1}^pa_kt^k$ 为多项式, 且 $Q(1)\ne0$.设 $E$ 为 Banach
空间, $T\in\LL(E)$.假设 $Q(T)\in\KK(E)$.
\begin{enumerate}[label=\arabic*., leftmargin=2.2em, itemsep=0.3em]
	\item 证明 $\dim N(I-T)<\infty$, 且 $R(I-T)$ 闭.更一般地, 证明对每个闭
	      子空间 $E_0\subset E$, $(I-T)(E_0)$ 闭.
	      [提示: 写出 $Q(1)-Q(t)=\tilde Q(t)(1-t)$, 其中 $\tilde Q$ 为某个
	      多项式, 并应用习题  \ref{exer:6.9}.]
	\item 证明 $N(I-T)=\{0\}\iff R(I-T)=E$.
	\item 证明 $\dim N(I-T)=\dim N(I-T^\star)$.
	      [第 2, 3 问的提示: 使用与定理  \ref{thm:6.6} 的证明相同的方法.]
\end{enumerate}
\end{exercise}

\begin{exercise}
\label{exer:6.11}
设 $K$ 为紧度量空间, $E=C(K;\R)$ 赋以通常的范数
$\norm{u}=\max_{x\in K}\abs{u(x)}$.设 $F\subset E$ 为闭子空间.假设每个函数
$u\in F$ 都是 Hölder 连续的, 即
\begin{equation*}
\forall u\in F\ \ \exists\alpha\in(0,1]\ \text{与}\ \exists L
\ \text{使得}\ \abs{u(x)-u(y)}\le Ld(x,y)^\alpha
\quad\forall x,y\in K.
\end{equation*}
本题的目的是证明 $F$ 是有限维的.
\begin{enumerate}[label=\arabic*., leftmargin=2.2em, itemsep=0.3em]
	\item 证明存在常数 $\gamma\in(0,1]$ 与 $C\ge0$ (二者都与 $u$ 无关),
	      使得
	      \begin{equation*}
	      \abs{u(x)-u(y)}\le C\norm{u}d(x,y)^\gamma
	      \quad\forall u\in F,\ \forall x,y\in K.
	      \end{equation*}
	      [提示: 应用 Baire 纲定理 (定理  \ref{thm:2.1}), 取
	      $F_n=\{u\in F;\ \abs{u(x)-u(y)}\le nd(x,y)^{1/n}
	      \ \forall x,y\in K\}$.]
	\item 证明 $B_F$ 紧, 并作出结论.
\end{enumerate}
\end{exercise}

\begin{exercise}
\label{exer:6.12}
\emph{J.-L.~Lions 的一个引理.}
设 $X$, $Y$, $Z$ 为三个 Banach 空间, 范数分别为 $\norm{\cdot}_X$, $\norm{\cdot}_Y$, $\norm{\cdot}_Z$.假设 $X\subset Y$ 的嵌入是紧的,
$Y\subset Z$ 的嵌入是连续的.证明
\begin{equation*}
\forall\varepsilon>0\ \ \exists C_\varepsilon\ge0\ \text{使得}\
\norm{u}_Y\le\varepsilon\norm{u}_X+C_\varepsilon\norm{u}_Z
\quad\forall u\in X.
\end{equation*}
[提示: 用反证法.]

\emph{应用.} 证明
\begin{equation*}
\forall\varepsilon>0\ \ \exists C_\varepsilon\ge0\ \text{使得}\
\max_{[0,1]}\abs{u}\le\varepsilon\max_{[0,1]}\abs{u'}+C_\varepsilon\norm{u}_{L^1}
\quad\forall u\in C^1([0,1]).
\end{equation*}
\end{exercise}

\begin{exercise}
\label{exer:6.13}
设 $E$ 与 $F$ 为两个 Banach 空间, 范数分别为 $\norm{\cdot}_E$ 与
$\norm{\cdot}_F$.假设 $E$ 自反.设 $T\in\KK(E,F)$.考虑 $E$ 上的另一个
范数 $\abs{\cdot}$, 它比 $\norm{\cdot}_E$ 弱, 即
$\abs{u}\le C\norm{u}_E\ \forall u\in E$.证明
\begin{equation*}
\forall\varepsilon>0\ \ \exists C_\varepsilon\ge0\ \text{使得}\
\norm{Tu}_F\le\varepsilon\norm{u}_E+C_\varepsilon\abs{u}
\quad\forall u\in E.
\end{equation*}
证明当 $E$ 不自反时结论可能不成立.
[提示: 取 $E=C([0,1])$, $F=\R$, $\norm{u}=\norm{u}_{L^\infty}$,
$\abs{u}=\norm{u}_{L^1}$.]
\end{exercise}

\begin{exercise}
\label{exer:6.14}
设 $E$ 为 Banach 空间, $T\in\LL(E)$ 且 $\norm{T}<1$.
\begin{enumerate}[label=\arabic*., leftmargin=2.2em, itemsep=0.3em]
	\item 证明 $(I-T)$ 是双射, 且
	      \begin{equation*}
	      \norm{(I-T)^{-1}}\le\frac{1}{1-\norm{T}}.
	      \end{equation*}
	\item 令 $S_n=I+T+\cdots+T^{n-1}$.证明
	      \begin{equation*}
	      \norm{S_n-(I-T)^{-1}}\le\frac{\norm{T}^n}{1-\norm{T}}.
	      \end{equation*}
\end{enumerate}
\end{exercise}

\begin{exercise}
\label{exer:6.15}
设 $E$ 为 Banach 空间, $T\in\LL(E)$.
\begin{enumerate}[label=\arabic*., leftmargin=2.2em, itemsep=0.3em]
	\item 设 $\lambda\in\R$ 满足 $\abs{\lambda}>\norm{T}$.证明
	      \begin{equation*}
	      \norm{I+\lambda(T-\lambda I)^{-1}}\le\frac{\norm{T}}{\abs{\lambda}-\norm{T}}.
	      \end{equation*}
	\item 设 $\lambda\in\rho(T)$.验证
	      \begin{equation*}
	      (T-\lambda I)^{-1}T=T(T-\lambda I)^{-1},
	      \end{equation*}
	      并证明
	      \begin{equation*}
	      \dist(\lambda,\sigma(T))\ge\frac{1}{\norm{(T-\lambda I)^{-1}}}.
	      \end{equation*}
	\item 假设 $0\in\rho(T)$.证明
	      \begin{equation*}
	      \sigma(T^{-1})=\frac{1}{\sigma(T)}.
	      \end{equation*}
\end{enumerate}

下文假设 $1\in\rho(T)$; 令
\begin{equation*}
U=(T+I)(T-I)^{-1}=(T-I)^{-1}(T+I).
\end{equation*}

\begin{enumerate}[label=\arabic*., leftmargin=2.2em, start=4, itemsep=0.3em]
	\item 验证 $1\in\rho(U)$, 并用 $T$ 给出 $(U-I)^{-1}$ 的简单表达式.
	\item 证明 $T=(U+I)(U-I)^{-1}$.
	\item 考虑函数 $f(t)=\frac{t+1}{t-1}$, $t\in\R$.证明
	      \begin{equation*}
	      \sigma(U)=f(\sigma(T)).
	      \end{equation*}
\end{enumerate}
\end{exercise}

\begin{exercise}
\label{exer:6.16}
设 $E$ 为 Banach 空间, $T\in\LL(E)$.
\begin{enumerate}[label=\arabic*., leftmargin=2.2em, itemsep=0.3em]
	\item 假设 $T^2=I$.证明 $\sigma(T)\subset\{-1,+1\}$, 并对
	      $\lambda\ne\pm1$ 确定 $(T-\lambda I)^{-1}$.
	\item 更一般地, 假设存在整数 $n\ge2$ 使得 $T^n=I$.证明
	      $\sigma(T)\subset\{-1,+1\}$, 并对 $\lambda\ne\pm1$ 确定
	      $(T-\lambda I)^{-1}$.
	\item 假设存在整数 $n\ge2$ 使得 $T^n=0$.证明 $\sigma(T)=\{0\}$, 并对
	      $\lambda\ne0$ 确定 $(T-\lambda I)^{-1}$.
	\item 假设存在整数 $n\ge2$ 使得 $\norm{T^n}<1$.证明 $I-T$ 是双射, 并用
	      $(I-T^n)^{-1}$ 与 $T$ 的迭代给出 $(I-T)^{-1}$ 的表达式.
\end{enumerate}
\end{exercise}

\begin{exercise}
\label{exer:6.17}
设 $E=\ell^p$, $1\le p\le\infty$, $(\lambda_n)$ 为 $\R$ 中的有界序列.
考虑由
\begin{equation*}
Mx=(\lambda_1x_1,\lambda_2x_2,\dots,\lambda_nx_n,\dots)
\end{equation*}
定义的乘法算子 $M\in\LL(E)$, 其中 $x=(x_1,x_2,\dots,x_n,\dots)$.确定
$EV(M)$ 与 $\sigma(M)$.
\end{exercise}

\begin{exercise}
\label{exer:6.18}
\emph{移位算子的谱性质.}
$E=\ell^2$ 中的元素记作 $x=(x_1,x_2,\dots,x_n,\dots)$.考虑算子
\begin{equation*}
S_rx=(0,x_1,x_2,\dots,x_{n-1},\dots),\qquad
S_\ell x=(x_2,x_3,x_4,\dots,x_{n+1},\dots),
\end{equation*}
它们分别称为\emph{右移位}与\emph{左移位}.
\begin{enumerate}[label=\arabic*., leftmargin=2.2em, itemsep=0.3em]
	\item 确定 $\norm{S_r}$ 与 $\norm{S_\ell}$.$S_r$ 或 $S_\ell$ 属于
	      $\KK(E)$ 吗?
	\item 证明 $EV(S_r)=\emptyset$.
	\item 证明 $\sigma(S_r)=[-1,+1]$.
	\item 证明 $EV(S_\ell)=(-1,+1)$.确定对应的特征空间.
	\item 证明 $\sigma(S_\ell)=[-1,+1]$.
	\item 确定 $S_r^\star$ 与 $S_\ell^\star$.
	\item 证明对每个 $\lambda\in(-1,+1)$, 空间 $R(S_r-\lambda I)$ 与
	      $R(S_\ell-\lambda I)$ 都闭.给出这些空间的显式表示.
	      [提示: 应用定理  \ref{thm:2.19} 与  \ref{thm:2.20}.]
	\item 证明空间 $R(S_r\pm I)$ 与 $R(S_\ell\pm I)$ 稠密但不闭.
\end{enumerate}

考虑由
\begin{equation*}
Mx=(\alpha_1x_1,\alpha_2x_2,\dots,\alpha_nx_n,\dots)
\end{equation*}
定义的乘法算子 $M$, 其中 $(\alpha_n)$ 是 $\R$ 中的有界序列.

\begin{enumerate}[label=\arabic*., leftmargin=2.2em, start=9, itemsep=0.3em]
	\item 确定 $EV(S_r\circ M)$.
	\item 假设当 $n\to\infty$ 时 $\alpha_n\to\alpha$.证明
	      \begin{equation*}
	      \sigma(S_r\circ M)=[-\abs{\alpha},+\abs{\alpha}].
	      \end{equation*}
	      [提示: 应用定理  \ref{thm:6.6}.]
	\item 假设对每个整数 $n$, $\alpha_{2n}=a$, $\alpha_{2n+1}=b$, 其中
	      $a\ne b$.确定 $\sigma(S_r\circ M)$.
	      [提示: 计算 $(S_r\circ M)^2$, 并应用习题  \ref{exer:6.16} 第 4 问.]
\end{enumerate}
\end{exercise}

\begin{exercise}
\label{exer:6.19}
设 $E$ 为 Banach 空间, $T\in\LL(E)$.
\begin{enumerate}[label=\arabic*., leftmargin=2.2em, itemsep=0.3em]
	\item 证明 $\sigma(T^\star)=\sigma(T)$.
	\item 举例说明 $EV(T)$ 与 $EV(T^\star)$ 之间一般没有包含关系.
	      [提示: 考虑右移位与左移位.]
\end{enumerate}
\end{exercise}

\begin{exercise}
\label{exer:6.20}
设 $E=L^p(0,1)$, $1\le p<\infty$.给定 $u\in E$, 令
\begin{equation*}
Tu(x)=\int_0^xu(t)\,dt.
\end{equation*}
\begin{enumerate}[label=\arabic*., leftmargin=2.2em, itemsep=0.3em]
	\item 证明 $T\in\KK(E)$.
	\item 确定 $EV(T)$ 与 $\sigma(T)$.
	\item 当 $\lambda\in\rho(T)$ 时, 给出 $(T-\lambda I)^{-1}$ 的显式公式.
	\item 确定 $T^\star$.
\end{enumerate}
\end{exercise}

\begin{exercise}
\label{exer:6.21}
设 $V$ 与 $H$ 为两个 Banach 空间, 范数分别为 $\norm{\cdot}$ 与
$\abs{\cdot}$, 满足
\begin{equation*}
V\subset H\ \text{的嵌入是紧的}.
\end{equation*}
设 $p(u)$ 为 $V$ 上的半范数, 使得 $p(u)+\abs{u}$ 是 $V$ 上与
$\norm{\cdot}$ 等价的范数.令
\begin{equation*}
N=\{u\in V;\ p(u)=0\},
\end{equation*}
并对 $u\in V$ 令 $\dist(u,N)=\inf_{v\in N}\norm{u-v}$.
\begin{enumerate}[label=\arabic*., leftmargin=2.2em, itemsep=0.3em]
	\item 证明 $N$ 是有限维空间.
	      [提示: 考虑 $N$ 中赋以范数 $\abs{\cdot}$ 的单位球.]
	\item 证明存在常数 $K_1>0$, 使得
	      \begin{equation*}
	      p(u)\le K_1\dist(u,N)\quad\forall u\in V.
	      \end{equation*}
	\item 证明存在常数 $K_2>0$, 使得
	      \begin{equation*}
	      K_2\dist(u,N)\le p(u)\quad\forall u\in V.
	      \end{equation*}
	      [提示: 用反证法.假设存在 $V$ 中的序列 $(u_n)$, 使得
	      $\dist(u_n,N)=1\ \forall n$ 且 $p(u_n)\to0$.]
\end{enumerate}
\end{exercise}

\begin{exercise}
\label{exer:6.22}
设 $E$ 为 Banach 空间, $T\in\LL(E)$.给定多项式
$Q(t)=\sum_{k=0}^pa_kt^k$ ($a_k\in\R$), 令
$Q(T)=\sum_{k=0}^pa_kT^k$.
\begin{enumerate}[label=\arabic*., leftmargin=2.2em, itemsep=0.3em]
	\item 证明 $Q(EV(T))\subset EV(Q(T))$.
	\item 证明 $Q(\sigma(T))\subset\sigma(Q(T))$.
	\item 在 $E=\R^2$ 中构造一个例子, 使上述包含是严格的.
\end{enumerate}

下文假设 $E$ 是 Hilbert 空间 (与其对偶空间 $H^\star$ 等同), 且
$T^\star=T$.

\begin{enumerate}[label=\arabic*., leftmargin=2.2em, start=4, itemsep=0.3em]
	\item 这里假设多项式 $Q$ 没有实根, 即
	      $Q(t)\ne0\ \forall t\in\R$.证明 $Q(T)$ 是双射.
	      [提示: 先考虑 $Q$ 是二次多项式的情形, 更确切地说,
	      $Q(t)=t^2+1$.]
	\item 由此推出对每个多项式 $Q$, 都有
	      \begin{enumerate}[label=(\roman*), leftmargin=2.4em, itemsep=0.1em]
	      	\item $Q(EV(T))=EV(Q(T))$,
	      	\item $Q(\sigma(T))=\sigma(Q(T))$.
	      \end{enumerate}
	      [提示: 写出
	      $Q(t)-\lambda=(t-t_1)(t-t_2)\cdots(t-t_q)\tilde Q(t)$, 其中
	      $t_1,t_2,\dots,t_q$ 是 $Q(t)-\lambda$ 的实根, $\tilde Q$ 无实根.]
\end{enumerate}
\end{exercise}

\begin{exercise}
\label{exer:6.23}
\emph{谱半径.}
设 $E$ 为 Banach 空间, $T\in\LL(E)$.令
\begin{equation*}
a_n=\log\norm{T^n},\quad n\ge1.
\end{equation*}
\begin{enumerate}[label=\arabic*., leftmargin=2.2em, itemsep=0.3em]
	\item 验证
	      \begin{equation*}
	      a_{i+j}\le a_i+a_j\quad\forall i,j\ge1.
	      \end{equation*}
	\item 由此推出
	      \begin{equation*}
	      \lim_{n\to+\infty}(a_n/n)\ \text{存在, 且等于}\ \inf_{m\ge1}(a_m/m).
	      \end{equation*}
	      [提示: 固定整数 $m\ge1$.对任意整数 $n\ge1$, 写出
	      $n=mq+r$, 其中 $q=[\frac nm]$ 是不超过 $n/m$ 的最大整数,
	      $0\le r<m$.注意 $a_n\le\frac nm a_m+a_r$.]
	\item 作出结论: $r(T)=\lim_{n\to\infty}\norm{T^n}^{1/n}$ 存在, 且
	      $r(T)\le\norm{T}$.在 $E=\R^2$ 中构造一个例子, 使得 $r(T)=0$
	      而 $\norm{T}=1$.
\end{enumerate}

数 $r(T)$ 称为 $T$ 的\emph{谱半径}.

\begin{enumerate}[label=\arabic*., leftmargin=2.2em, start=4, itemsep=0.3em]
	\item 证明 $\sigma(T)\subset[-r(T),+r(T)]$.由此推出: 若
	      $\sigma(T)\ne\emptyset$, 则
	      $\max\{\abs{\lambda};\ \lambda\in\sigma(T)\}\le r(T)$.
	      [提示: 注意若 $\lambda\in\sigma(T)$, 则
	      $\lambda^n\in\sigma(T^n)$; 见习题  \ref{exer:6.22}.]
	\item 在 $E=\R^3$ 中构造一个例子, 使得 $\sigma(T)=\{0\}$, 而
	      $r(T)=1$.
\end{enumerate}

下文取 $E=L^p(0,1)$, $1\le p\le\infty$.考虑由
\begin{equation*}
Tu(t)=\int_0^tu(s)\,ds
\end{equation*}
定义的算子 $T\in\LL(E)$.

\begin{enumerate}[label=\arabic*., leftmargin=2.2em, start=6, itemsep=0.3em]
	\item 用归纳法证明: 当 $n\ge2$ 时,
	      \begin{equation*}
	      (T^nu)(t)=\frac{1}{(n-1)!}\int_0^t(t-\tau)^{n-1}u(\tau)\,d\tau.
	      \end{equation*}
	\item 由此推出 $\norm{T^n}\le\frac{1}{n!}$.
	      [提示: 使用卷积积的一个不等式.]
	\item 证明 $T$ 的谱半径为 $0$.
	      [提示: 使用 Stirling 公式.]
	\item 证明 $\sigma(T)=\{0\}$.与习题  \ref{exer:6.20} 作比较.
\end{enumerate}
\end{exercise}

\begin{exercise}
\label{exer:6.24}
假设 $T\in\LL(H)$ 自伴.
\begin{enumerate}[label=\arabic*., leftmargin=2.2em, itemsep=0.3em]
	\item 证明下述性质等价:
	      \begin{enumerate}[label=(\roman*), leftmargin=2.4em, itemsep=0.1em]
	      	\item $(Tu,u)\ge0\ \forall u\in H$,
	      	\item $\sigma(T)\subset[0,\infty)$.
	      \end{enumerate}
	      [提示: 应用命题  \ref{prop:6.9}.]
	\item 证明下述性质等价:
	      \begin{enumerate}[label=(\roman*), leftmargin=2.4em, start=3, itemsep=0.1em]
	      	\item $\norm{T}\le1$ 且 $(Tu,u)\ge0\ \forall u\in H$,
	      	\item $0\le(Tu,u)\le\abs{u}^2\ \forall u\in H$,
	      	\item $\sigma(T)\subset[0,1]$,
	      	\item $(Tu,u)\ge\abs{Tu}^2\ \forall u\in H$.
	      \end{enumerate}
	      [提示: 为证 (v)$\Rightarrow$(vi), 对 $(T+\varepsilon I)^{-1}$
	      ($\varepsilon>0$) 应用命题  \ref{prop:6.9}.]
	\item 证明下述性质等价:
	      \begin{enumerate}[label=(\roman*), leftmargin=2.4em, start=7, itemsep=0.1em]
	      	\item $(Tu,u)\le\abs{Tu}^2\ \forall u\in H$,
	      	\item $(0,1)\subset\rho(T)$.
	      \end{enumerate}
	      [提示: 引入 $U=2T-I$.]
\end{enumerate}
\end{exercise}

\begin{exercise}
\label{exer:6.25}
设 $E$ 为 Banach 空间, $K\in\KK(E)$.证明存在
$M\in\LL(E)$, $\tilde M\in\LL(E)$ 以及有限秩投影算子 $P$, $\tilde P$,
使得
\begin{enumerate}[label=(\roman*), leftmargin=2.6em, itemsep=0.15em]
	\item $M\circ(I+K)=I-P$,
	\item $(I+K)\circ\tilde M=I-\tilde P$.
\end{enumerate}
[提示: 设 $X$ 为 $N(I+K)$ 在 $E$ 中的补空间.则 $(I+K)|_X$ 是由 $X$ 到
$R(I+K)$ 的双射.记其逆为 $M$.设 $Q$ 为 $E$ 到 $R(I+K)$ 的投影算子,
并令 $\tilde M=M\circ Q$.证明 (i) 与 (ii) 成立.]
\end{exercise}

\begin{exercise}
\label{exer:6.26}
\emph{从 Brouwer 不动点定理到 Schauder 不动点定理.}
本题假设下述结果已知 (证明见例如 K.~Deimling \bref{Deimling}{1}, A.~Granas--J.~Dugundji \bref{GranasDugundji}{1}, 或
L.~Nirenberg \bref{Nirenberg}{2}).

\smallskip
\noindent\emph{定理 (Brouwer).} 设 $F$ 为有限维空间, $Q\subset F$ 为非空紧凸集.
设 $f:Q\to Q$ 为连续映射.则 $f$ 有不动点, 即存在 $p\in Q$ 使得
$f(p)=p$.

\smallskip
我们的目标是证明下述结果.

\smallskip
\noindent\emph{定理 (Schauder).} 设 $E$ 为 Banach 空间, $C$ 为 $E$ 中的非空
闭凸集.设 $F:C\to C$ 为连续映射, 且 $F(C)\subset K$, 其中 $K$ 是 $C$ 的紧
子集.则 $F$ 在 $K$ 中有不动点.

\begin{enumerate}[label=\arabic*., leftmargin=2.2em, itemsep=0.3em]
	\item 给定 $\varepsilon>0$, 考虑 $K$ 的有限覆盖, 即
	      $K\subset\bigcup_{i\in I}B(y_i,\varepsilon/2)$, 其中 $I$ 有限,
	      $y_i\in K\ \forall i\in I$.定义函数
	      $q(x)=\sum_{i\in I}q_i(x)$, 其中
	      \begin{equation*}
	      q_i(x)=\max\{\varepsilon-\norm{Fx-y_i},0\}.
	      \end{equation*}
	      验证 $q$ 在 $C$ 上连续, 且 $q(x)\ge\varepsilon/2\ \forall x\in C$.
	\item 令
	      \begin{equation*}
	      F_\varepsilon(x)=\frac{\sum_{i\in I}q_i(x)y_i}{q(x)},\quad x\in C.
	      \end{equation*}
	      证明 $F_\varepsilon:C\to C$ 连续, 且
	      \begin{equation*}
	      \norm{F_\varepsilon(x)-F(x)}\le\varepsilon,\quad\forall x\in C.
	      \end{equation*}
	\item 证明 $F_\varepsilon$ 有不动点 $x_\varepsilon\in C$.
	      [提示: 令 $Q=\operatorname{conv}\bigl(\bigcup_{i\in I}\{y_i\}\bigr)$.
	      验证 $F_\varepsilon|_Q$ 在 $Q$ 中有不动点 $x_\varepsilon\in Q$.]
	\item 证明对某个序列 $\varepsilon_n\to0$, $(x_{\varepsilon_n})$ 收敛于
	      某个极限 $x\in C$.证明 $F(x)=x$.
\end{enumerate}
\end{exercise}

